\studySubsection{index-methods}{方法索引}

\begin{knowledgeBox}
\knowledgeIndexAnchor[MATH1-KN-CALC-0059]{七种未定式的统一化归}
\textbf{七种未定式的统一化归}

\begin{itemize}
  \item \textbf{识别范围：} \(\frac00\)、\(\frac{\infty}{\infty}\)、\(0\cdot\infty\)、\(\infty-\infty\)、\(\infty^0\)、\(0^0\)、\(1^\infty\)。
  \item \textbf{统一主线：} 两种商式直接计算；积式通过取倒数化为商式；差式通过通分、有理化或合并化为商式；三种幂式令 \(y=f(x)^{g(x)}\)，先求 \(\lim g(x)\ln f(x)\)，再取指数。
  \item \textbf{工具优先级：} 先约分、提取非零因子、有理化、通分、等价无穷小或同除最高阶量，再按条件使用泰勒展开或洛必达法则。
  \item \textbf{高阶相消：} 等价无穷小通常只在乘除结构中直接替换；加减主部相消时必须展开到第一个不为零的高阶项，例如 \(x-\ln(1+x)\sim\frac12x^2\)。
  \item \textbf{参数与特殊结构：} 抓最高阶或最低阶项前检查对应系数是否可能为零；含取整、振荡或有界扰动时，优先隔离有界误差并用夹逼。
  \item \textbf{已知极限反求未知：} 将 \(F(x)\to L\) 改写为 \(F(x)=L+\alpha(x)\)、\(\alpha(x)\to0\)，再从题设关系中解出未知函数；高阶项由泰勒展开确定。
  \item \textbf{使用提醒：} 洛必达法则只直接处理 \(\frac00\) 与 \(\frac{\infty}{\infty}\)；幂式取对数前必须检查底数最终为正；每次变形后都要重新判型。
  \item \textbf{关联知识点：} \knowledgeRef[MATH1-KN-CALC-0059]{七种未定式的统一化归}（第 1 讲“函数极限与连续”）。
\end{itemize}

\knowledgeIndexAnchor[MATH1-KN-CALC-0002]{整体换元}
\legacyKnowledgeIndexAnchor{substitution-as-whole}
\textbf{整体换元反求函数表达式}

\begin{itemize}
  \item \textbf{适用题型：} 已知 \(f(\varphi(u))\)，要求 \(f(x)\) 或 \(f(t)\)。
  \item \textbf{触发特征：} 题设中出现 \(u+\frac{1}{u}\)、\(u-\frac{1}{u}\)、\(u^2+\frac{1}{u^2}\) 等可互相转化的结构。
  \item \textbf{核心步骤：} 先将原参数重命名，再把右端化成只含 \(\varphi(u)\) 的表达式，最后令 \(t=\varphi(u)\)。
  \item \textbf{关联题目：} \problemRef{MATH1-CALC-0001}。
\end{itemize}

\textbf{反解参数后消元}

\begin{itemize}
  \item \textbf{适用题型：} 已知 \(f(\varphi(u))\)，且 \(\varphi(u)=t\) 可以转化为低次代数方程。
  \item \textbf{触发特征：} 内层变量如 \(t=u+\frac{1}{u}\) 可化为 \(u^2-tu+1=0\)，从而用 \(1+u^2=tu\) 等关系消去 \(u\)。
  \item \textbf{核心步骤：} 先设 \(t=\varphi(u)\)，建立 \(u\) 与 \(t\) 的方程关系，再把右端分子分母分别化成含 \(t\) 的倍数，最后约去参数。
  \item \textbf{使用提醒：} 本法常比整体换元繁琐；也可显式解出 \(u\) 的分支后代回，但若不同分支会导致右端不同值，必须额外检查分支是否影响函数定义。
  \item \textbf{关联题目：} \problemRef{MATH1-CALC-0001}。
\end{itemize}

\knowledgeIndexAnchor[MATH1-KN-CALC-0008]{反函数}
\legacyKnowledgeIndexAnchor{inverse-function}
\textbf{求反函数五步法}

\begin{itemize}
  \item \textbf{适用题型：} 已知 \(y=f(x)\)，要求写出 \(f^{-1}(x)\) 或判断反函数的定义域、值域。
  \item \textbf{触发特征：} 题目要求“求反函数”“写出反函数表达式”，或需要从函数值唯一倒推自变量。
  \item \textbf{核心步骤：} 明确原定义域；由 \(y=f(x)\) 解出 \(x\)；检查唯一性；交换 \(x,y\)；写明反函数定义域为原函数值域、反函数值域为原函数定义域。
  \item \textbf{使用提醒：} \(f^{-1}(x)\) 表示反函数，不是 \(\frac{1}{f(x)}\)；若解出多个分支，必须限制定义域或说明分支选择。
  \item \textbf{关联知识点：} \knowledgeRef[MATH1-KN-CALC-0008]{反函数}（第 1 讲“反函数的定义、存在条件与图像关系”）。
\end{itemize}

\textbf{水平线判别法}

\begin{itemize}
  \item \textbf{适用题型：} 判断一个已知函数是否存在反函数，尤其适合图像题和分段函数题。
  \item \textbf{触发特征：} 需要判断同一个函数值是否可能对应多个自变量。
  \item \textbf{核心步骤：} 在确认图像已经表示函数后，作水平线；若任意水平线与图像至多一个交点，则函数一一对应，有反函数。
  \item \textbf{使用提醒：} 铅直线判定“是不是函数”，水平线判定“有没有反函数”，两者不能混用。
  \item \textbf{关联知识点：} \knowledgeRef[MATH1-KN-CALC-0008]{反函数}（第 1 讲“反函数的定义、存在条件与图像关系”）。
\end{itemize}

\textbf{对数反函数中的倒数有理化法}

\begin{itemize}
  \item \textbf{适用题型：} 求 \(y=\ln\left(x+\sqrt{x^2+1}\right)\) 或同型函数的反函数。
  \item \textbf{触发特征：} 对数真数含有 \(x+\sqrt{x^2+1}\)，其倒数可通过有理化化为 \(\sqrt{x^2+1}-x\)。
  \item \textbf{核心步骤：} 由 \(e^y=x+\sqrt{x^2+1}\) 得一式；再由倒数有理化得 \(e^{-y}=\sqrt{x^2+1}-x\)；把两式看成关于 \(x\) 与 \(\sqrt{x^2+1}\) 的方程组，相减消去根式，得到 \(x=\frac12(e^y-e^{-y})\)。
  \item \textbf{思维本质：} 主动制造第二个方程，再联立消元；这与 \(f(x),f(1/x)\) 类题中的对偶代换方程组思维相通。
  \item \textbf{使用提醒：} 不要把 \(f^{-1}(x)\) 理解为倒数；求完表达式后必须写反函数定义域，即原函数值域。
  \item \textbf{关联题目：} \problemRef{MATH1-CALC-0003}。
\end{itemize}

\knowledgeIndexAnchor[MATH1-KN-CALC-0012]{复合函数}
\legacyKnowledgeIndexAnchor{composite-function}
\textbf{整体视角反求中间函数}

\begin{itemize}
  \item \textbf{适用题型：} 已知 \(f[\varphi(x)]\) 和外层函数 \(f\)，要求 \(\varphi(x)\) 及其定义域、值域。
  \item \textbf{触发特征：} 题设给出复合结果，但中间函数未知，例如 \(f(x)=x^2\)，\(f[\varphi(x)]=G(x)\)。
  \item \textbf{核心步骤：} 把 \(\varphi(x)\) 当成整体变量；由 \(f(t)\) 的表达式建立关于 \(\varphi(x)\) 的方程；再用附加条件确定分支，并检查定义域和值域。
  \item \textbf{使用提醒：} 若外层函数不是一一对应函数，例如平方函数，反求中间函数时必须处理正负分支。
  \item \textbf{关联题目：} \problemRef{MATH1-CALC-0004}。
\end{itemize}

\textbf{分段复合先判内层落段}

\begin{itemize}
  \item \textbf{适用题型：} 外层或内层为分段函数的复合函数，如 \(g[f(x)]\)。
  \item \textbf{触发特征：} 外层函数的分段条件依赖其输入，而输入本身是另一个函数值。
  \item \textbf{核心步骤：} 先写外层分段规则 \(g(u)\)，再令 \(u=f(x)\)，判断 \(f(x)\) 落入外层哪一段，最后代回化简。
  \item \textbf{使用提醒：} 外层函数的分段依据是内层输出，不是原来的 \(x\)；边界点要先按内层函数定义判断。
  \item \textbf{关联题目：} \problemRef{MATH1-CALC-0005}。
\end{itemize}

\knowledgeIndexAnchor[MATH1-KN-CALC-0029]{隐函数}
\legacyKnowledgeIndexAnchor{implicit-function}
\textbf{隐函数特殊点观察法}

\begin{itemize}
  \item \textbf{适用题型：} 已知 \(F(x,y)=0\) 确定隐函数，只要求 \(y(x_0)\)。
  \item \textbf{触发特征：} 代入 \(x=x_0\) 后关于 \(y\) 的方程容易观察出解。
  \item \textbf{核心步骤：} 先代入 \(x_0\)，再观察明显根；若需严谨，补充单调性或图像交点唯一性。
  \item \textbf{使用提醒：} 不要为了求一个点强行显化整个隐函数。
  \item \textbf{关联知识点：} \knowledgeRef[MATH1-KN-CALC-0029]{隐函数}（第 1 讲“隐函数的概念与观察法”）。
\end{itemize}

\knowledgeIndexAnchor[MATH1-KN-CALC-0017]{有界性}
\legacyKnowledgeIndexAnchor{boundedness}
\textbf{有界性放缩找上界}

\begin{itemize}
  \item \textbf{适用题型：} 证明函数在某区间上有界。
  \item \textbf{触发特征：} 目标可转化为寻找 \(M\)，使 \(|f(x)|\le M\)。
  \item \textbf{核心步骤：} 先取绝对值，再用基本不等式、配方、分母放缩或端点分析寻找统一上界。
  \item \textbf{方法价值：} 若只需证明有界，放缩通常能替代“求导判断单调性、找极值点、比较最大值”的较重路线。
  \item \textbf{使用提醒：} 讨论有界性必须说明区间；若除以含 \(x\) 的量，要单独检查零点；若题目要求最值和取等条件，再考虑极值法。
  \item \textbf{关联题目：} \problemRef{MATH1-CALC-0006}。
\end{itemize}

\knowledgeIndexAnchor[MATH1-KN-CALC-0021]{单调性}
\legacyKnowledgeIndexAnchor{monotonicity}
\textbf{单调性定义法追踪不等号}

\begin{itemize}
  \item \textbf{适用题型：} 判断图像变换、复合或变号后的单调性。
  \item \textbf{触发特征：} 题设给出单调性定义或乘积判别式，而不是导数。
  \item \textbf{核心步骤：} 设新函数 \(F(x)\)，任取 \(x_1<x_2\)，追踪输入变换和外层变号造成的不等号变化。
  \item \textbf{使用提醒：} \(f(-x)\) 会反转单调性，\(-f(-x)\) 会再反一次。
  \item \textbf{关联题目：} \problemRef{MATH1-CALC-0007}。
\end{itemize}

\knowledgeIndexAnchor[MATH1-KN-CALC-0024]{奇函数}
\legacyKnowledgeIndexAnchor{odd-function}
\textbf{函数方程赋值法}

\begin{itemize}
  \item \textbf{适用题型：} 已知对任意变量成立的函数方程，证明奇偶性或求特殊值。
  \item \textbf{触发特征：} 题设含“任意 \(x,y\)”并出现 \(f(x+y)\)、\(f(x-y)\) 等结构。
  \item \textbf{核心步骤：} 先令变量为 \(0\) 求 \(f(0)\)，再令两个变量互为相反数或相等，制造目标关系。
  \item \textbf{使用提醒：} 证明奇函数时，目标是 \(f(-x)=-f(x)\)。
  \item \textbf{关联题目：} \problemRef{MATH1-CALC-0008}。
\end{itemize}

\knowledgeIndexAnchor[MATH1-KN-CALC-0026]{周期函数}
\legacyKnowledgeIndexAnchor{periodic-function}
\textbf{递推式证明周期性}

\begin{itemize}
  \item \textbf{适用题型：} 已知 \(f(x)=f(x-a)+\varphi(x)\)，证明某个 \(T\) 是周期。
  \item \textbf{触发特征：} 题设能把 \(f(x+a)\)、\(f(x+2a)\) 逐步拉回 \(f(x)\)。
  \item \textbf{核心步骤：} 从目标 \(f(x+T)\) 出发，多次套用题设，再利用新增项相互抵消。
  \item \textbf{使用提醒：} 证明 \(T\) 是周期不等于证明最小正周期。
  \item \textbf{关联题目：} \problemRef{MATH1-CALC-0009}。
\end{itemize}

\knowledgeIndexAnchor[MATH1-KN-CALC-0005]{倒数代换}
\legacyKnowledgeIndexAnchor{reciprocal-substitution}
\textbf{对偶代换建立方程组}

\begin{itemize}
  \item \textbf{适用题型：} 关系式中同时出现 \(f(x)\) 与 \(f\left(\frac{1}{x}\right)\)、\(f(-x)\)、\(f(a-x)\) 等成对函数值。
  \item \textbf{触发特征：} 目标是求 \(f(x)\)，但题设中还混有另一个函数值，需要通过自变量替换得到第二个方程。
  \item \textbf{核心步骤：} 先把混合出现的函数值看作未知量，再做对应代换，如 \(x\mapsto \frac{1}{x}\)，联立方程并消去非目标项。
  \item \textbf{使用提醒：} 代换前必须检查定义域封闭性；例如本题定义域为 \((0,+\infty)\)，所以 \(x>0\) 时 \(\frac{1}{x}>0\)。
  \item \textbf{关联题目：} \problemRef{MATH1-CALC-0002}。
\end{itemize}
\end{knowledgeBox}

% BEGIN calculus-foundation:method-index
\begin{knowledgeBox}
\knowledgeIndexAnchor[MATH1-KN-CALC-0310]{定义域与复合函数}
\textbf{定义域与复合函数}

\begin{itemize}
  \item \textbf{识别：} “有意义”“定义域为”“已知f的定义域，求f(g(x))”。
  \item \textbf{流程：} ①列内层自身约束；②列内层值落入外层定义域；③取交集；④含参数时按临界点分段；⑤检查集合非空。
  \item \textbf{正文：} \knowledgeRef[MATH1-KN-CALC-0310]{定义域与复合函数}。
\end{itemize}

\knowledgeIndexAnchor[MATH1-KN-CALC-0311]{分段点连续与可导参数}
\textbf{分段点连续与可导参数}

\begin{itemize}
  \item \textbf{识别：} “在x=0处连续/可导”“确定a,b”。
  \item \textbf{流程：} ①算左右极限和点值，先解连续条件；②算左右导数，解可导条件；③二阶问题再逐侧求导；④回代核验。
  \item \textbf{正文：} \knowledgeRef[MATH1-KN-CALC-0311]{分段点连续与可导参数}。
\end{itemize}

\knowledgeIndexAnchor[MATH1-KN-CALC-0312]{常规极限方法选择}
\textbf{常规极限方法选择}

\begin{itemize}
  \item \textbf{识别：} 代入后出现未定式；含基本小量、根式、积分、幂指或振荡因子。
  \item \textbf{流程：} ①先判趋近过程与型态；②代数化简；③乘商看等价无穷小；④高阶消去用Taylor；⑤商型可比较洛必达；⑥振荡问题用夹逼或见证数列；⑦最后核验条件。
  \item \textbf{振荡见证法：} 对“发散振幅 \(\times\) 振荡因子”，找一列峰值点证明无界，再找一列零点或小值点否定无穷大量；本质是区分“邻域中存在大值”与“邻域内所有值最终都大”。
  \item \textbf{正文：} \knowledgeRef[MATH1-KN-CALC-0312]{常规极限方法选择}。
\end{itemize}

\knowledgeIndexAnchor[MATH1-KN-CALC-0313]{幂指型极限}
\textbf{幂指型极限}

\begin{itemize}
  \item \textbf{识别：} “$f(x)^{g(x)}$”“$x^x$”“底数趋1”。
  \item \textbf{流程：} ①确认底数在邻域为正；②令y为原式并取ln；③求 $g\ln f$ 的极限；④用等价/Taylor/洛必达；⑤指数还原；⑥处理±∞。
  \item \textbf{正文：} \knowledgeRef[MATH1-KN-CALC-0313]{幂指型极限}。
\end{itemize}

\knowledgeIndexAnchor[MATH1-KN-CALC-0314]{递推数列极限}
\textbf{递推数列极限}

\begin{itemize}
  \item \textbf{识别：} “递推定义”“证明数列收敛”“求lim a\_n”。
  \item \textbf{流程：} ①解不动点方程找候选；②由初值构造不变区间；③证明单调或压缩；④由单调有界得收敛；⑤代入连续递推求极限；⑥排除其他根。
  \item \textbf{正文：} \knowledgeRef[MATH1-KN-CALC-0314]{递推数列极限}。
\end{itemize}
\end{knowledgeBox}
% END calculus-foundation:method-index
